% Glossary entries — extend as needed. Used through \gls{key} in the chapters.

\newglossaryentry{ats}{
	name={ATS},
	description={Applicant Tracking System --- enterprise software used by HR teams to receive, store and manage candidate applications across recruitment processes}
}

\newglossaryentry{cefr}{
	name={CEFR},
	description={Common European Framework of Reference for Languages --- a six-level scale (A1 to C2) defined by the Council of Europe to describe language proficiency}
}

\newglossaryentry{gdpr}{
	name={GDPR},
	description={General Data Protection Regulation (EU) 2016/679 --- European regulation governing the processing of personal data of individuals in the European Union}
}

\newglossaryentry{llm}{
	name={LLM},
	description={Large Language Model --- a neural network trained on broad text corpora capable of producing and analysing natural language at scale}
}

\newglossaryentry{rls}{
	name={RLS},
	description={Row-Level Security --- a PostgreSQL feature that lets the database evaluate per-row access policies on every query}
}

\newglossaryentry{onion}{
	name={Onion architecture},
	description={Software architecture style organising code in concentric layers (Domain, Application, Infrastructure, Presentation), with dependencies pointing only inward toward the Domain}
}

\newglossaryentry{po}{
	name={Product Owner},
	description={Scrum role accountable for maximising the value of the product, owning and ordering the product backlog and acting as the main interface between development team and stakeholders}
}

\newglossaryentry{scrum}{
	name={Scrum},
	description={Agile framework for product development based on time-boxed iterations (sprints), explicit roles (Product Owner, Scrum Master, Developers) and a small set of artefacts and events}
}

\newglossaryentry{blended}{
	name={BlendEd},
	description={International academic--industry partnership programme in which multinational student teams from European universities research, design and implement solutions for problems presented by partner companies. In this edition one of the partner companies was adidas Global Business Services Porto}
}

\newglossaryentry{vercel}{
	name={Vercel},
	description={Cloud platform specialising in the deployment of front-end and full-stack web applications, with native support for Next.js, serverless API routes, edge functions and automatic CI/CD from a Git repository}
}

\newglossaryentry{oauth}{
	name={OAuth},
	description={Open authorisation standard (OAuth 2.0) that allows a third-party application to obtain delegated access to a user's resources at another service (e.g.\ Google) without exposing credentials}
}

\newglossaryentry{jwt}{
	name={JWT},
	description={JSON Web Token --- a compact, URL-safe token format defined in RFC 7519 for encoding claims between parties; typically signed (JWS) or encrypted (JWE) and used to carry session or identity assertions}
}

\newglossaryentry{hmac}{
	name={HMAC},
	description={Hash-based Message Authentication Code --- a cryptographic construction (defined in RFC 2104) that combines a secret key with a hash function to produce a tag that verifies both the integrity and the authenticity of a message}
}

\newglossaryentry{rest}{
	name={REST},
	description={Representational State Transfer --- an architectural style for distributed hypermedia systems that uses standard HTTP verbs (GET, POST, PUT, PATCH, DELETE) and stateless, resource-oriented URIs}
}

\newglossaryentry{supabase}{
	name={Supabase},
	description={Open-source backend-as-a-service platform that bundles a managed PostgreSQL database, an authentication service (Supabase Auth), and object storage under a single API and dashboard. Used in TalentHub for data persistence, Google OAuth identity management, and CV file storage}
}

\newglossaryentry{typescript}{
	name={TypeScript},
	description={Statically typed superset of JavaScript developed by Microsoft that adds an optional type system and compiles to plain JavaScript; used in TalentHub for the entire Next.js application layer}
}

\newglossaryentry{cicd}{
	name={CI/CD},
	description={Continuous Integration / Continuous Delivery --- a software-delivery practice in which every code change is automatically built, type-checked and tested (CI) and, on success, deployed to a target environment (CD). TalentHub uses GitHub Actions for CI and Vercel for CD}
}

\newglossaryentry{hallucination}{
	name={Hallucination},
	description={In the context of Large Language Models, the generation of plausible-sounding but factually incorrect or entirely invented content. In structured-output tasks such as CV parsing, hallucination manifests as invented field values (e.g.\ a skill or date that does not appear in the source document)}
}

\newglossaryentry{schemadrift}{
	name={Schema drift},
	description={A failure mode of LLM-generated structured output in which the model returns a subtly malformed object --- missing required fields, using unexpected key names, or nesting values incorrectly --- that passes a surface-level check but violates the expected schema and corrupts downstream processing}
}

\newglossaryentry{moscow}{
	name={MoSCoW},
	description={A requirements prioritisation technique in which each requirement is labelled Must have, Should have, Could have, or Won't have (this time). The capitalised letters form the acronym; the lower-case letters are included to make it pronounceable}
}

\newglossaryentry{orm}{
	name={ORM},
	description={Object-Relational Mapper --- a software layer that translates between the object model of an application and the relational model of a database, generating SQL from method calls and mapping query results to typed objects. TalentHub omits an ORM in favour of direct queries through the Supabase client, preserving full control and visibility over the generated SQL}
}

\newglossaryentry{isoiec25010}{
	name={ISO/IEC~25010},
	description={International standard defining a product-quality model for software and software-intensive systems. It organises quality into eight top-level characteristics (functional suitability, performance efficiency, compatibility, usability, reliability, security, maintainability, portability), each further decomposed into sub-characteristics that guide both design decisions and verification criteria}
}

\newglossaryentry{whisper}{
	name={Whisper},
	description={Open-source automatic speech-recognition model released by OpenAI, able to transcribe speech to text across many languages. Considered in TalentHub as an offline fallback for the browser-based speech-to-text used by the AI Interviewer}
}

\newglossaryentry{zod}{
	name={Zod},
	description={TypeScript-first schema-declaration and validation library. In TalentHub it defines the expected shape of every use-case input and of the structured JSON returned by the LLMs, rejecting malformed data at the untrusted boundary}
}

\newglossaryentry{shadcn}{
	name={shadcn/ui},
	description={A collection of accessible React user-interface components built on Radix primitives and styled with Tailwind CSS. Rather than a dependency, the components are copied into the project and owned directly. Used for TalentHub's dialogs, tables, forms and other interface elements}
}

\newglossaryentry{tailwind}{
	name={Tailwind CSS},
	description={A utility-first CSS framework in which styling is composed from small, single-purpose classes applied directly in the markup, rather than from hand-written stylesheets. Used throughout the TalentHub front end}
}

\newglossaryentry{cookies}{
	name={Cookies},
	description={Small pieces of data that a server asks a browser to store and return on subsequent requests. In TalentHub, HTTP-only cookies carry the Supabase session so that authentication state persists across page loads without exposing tokens to client-side scripts}
}

\newglossaryentry{cors}{
	name={CORS},
	description={Cross-Origin Resource Sharing --- a browser security mechanism that uses HTTP headers to control whether a web page served from one origin may make requests to a different origin. Relevant to the calls between the Next.js application and the AI Interviewer service}
}

\newglossaryentry{mime}{
	name={MIME type},
	description={Multipurpose Internet Mail Extensions type --- a standardised label (e.g.\ \texttt{application/pdf}) that identifies the format of a file or payload. Used when validating uploaded CV files before they are stored and parsed}
}

\newglossaryentry{vitest}{
	name={Vitest},
	description={A fast unit-testing framework for JavaScript and TypeScript projects, sharing configuration with the Vite build tool. TalentHub uses it for its automated test suite (527 tests across 33 files at the time of writing)}
}

\newglossaryentry{v8}{
	name={V8},
	description={Google's high-performance JavaScript and WebAssembly engine, which powers Chrome and Node.js. In TalentHub it also underpins the \texttt{@vitest/coverage-v8} provider that measures test coverage}
}

\newglossaryentry{turbopack}{
	name={Turbopack},
	description={An incremental bundler for JavaScript and TypeScript, written in Rust and integrated with Next.js, used to build and serve the TalentHub application faster than the previous Webpack-based toolchain}
}

\newglossaryentry{gin}{
	name={GIN index},
	description={Generalised Inverted Index --- a PostgreSQL index type suited to columns holding composite values such as arrays or JSONB. In TalentHub a GIN index accelerates membership queries over the \texttt{fields\_of\_work} array on candidate experiences}
}
